\scanpage{11}
Pour la première qui est un tor. compact, la composante à l'infini est de tte façon nulle. Donc on connaît $\xi$ quand on connaît tous les $\xi_\ell$ ($\ell$ fini). Donc ici on \deleted{\unsure{a}} connaît $\xi$.

\uline{NB} On sait \deleted{\unsure{prouve}} que $\mathcal M$ \uline{n'a pas de foncteur fibre sur} $R$, \inserted{(donc \deleted{\unsure{pas pleinement}} $\mathcal M_R$ \deleted{\unsure{pas}} \unsure{non neutre})} soit à cause du fait que $\mathcal M$ est polarisable et que $\unsure{U|H}$ est disconnexe, soit parce que l'on \deleted{\unsure{aurait}} \inserted{a} une représentation de $G^0$ de degré $1$, rang $1$ (qui se prolonge à $G$) \deleted{$\unsure{=G\times\widehat{\mathbf Z}}$} \deleted{l'obstruction est dans $\operatorname{Ext}^2(\mathbf Z$}, qui donnerait un objet de \deleted{degré} rang $1$ et de degré $1$ (s'il y avait foncteur fibre), soit par l'argument de Serre sur la courbe elliptique d'inv. de Hasse nul.

A fortiori, $\xi_\infty\ne0$, ce qu'on peut écrire aussi \unsure{i.e.}, la composante de $\xi_\infty$ suivant $H^2(R,\mu_2)$ \unsure{est dite|est déjà} (l'image \unsure{de|dite} \inserted{dans $\unsure{H^2}$} $\xi$) $\ne0$; cela définit donc \inserted{la composante de} $\xi_p$ \deleted{\unsure{comme}} dans $H^2(Q_p,\mu_2)$ comme \inserted{l'unique} élément non nul \unsure{dudit}.
\mspage{10}
